When I first found the gap in reality, it was entirely by accident. For days after hearing a loud crash in the middle of the night, I'd set up a simple drop trap using a clean bedsheet and waited nightly for any sound until sleep caught me. After seven days, the tell-tale screech of surprise echoed against tile, and I threw off my covers, barrelling down the stairs before my sheets hit the mattress and wrenching the kitchen door wide open. Something was writhing beneath the sheet, but all I could make out was the base of its long armoured back, like a crocodile but with burnished red scales. I should have known better than to approach without caution. It must have smelled or heard my approach, because it released a hiss like a steam engine, stilling, before its powerful tail slammed my breath from my chest and my body into the kitchen cabinets behind me with a deafening crash. To be accurate, it knocked me into last week.

You would be correct to question how I arrived at my conclusion so quickly about where – or rather when – I was when I landed. In my confused state, it took me some seconds to be aware of anything other than the aching of my body, and some more to realise that the strange creature which had attacked me had disappeared. As had the sheet, any sign of entry, and the various cameras which I had set up around the kitchen in preparation. Leaning heavily on the counter with one arm and holding my head with the other, I would have concluded that the entire thing must have been some sort of hallucination, but a too-familiar voice called from the top of the stairs.

``I'll go check what it is,'' I'd loudly informed my housemate a week earlier, who’d been blearily peering out from her bedroom doorway. ``I'm sure it's just a pest or something.''

The me in the kitchen froze, certain that I hadn't found myself in the kitchen a week ago, and tried to observe an escape from the situation that was about to occur. I couldn't remain, unless I wanted to risk creating an offshoot timeline and potentially dooming myself (both of us!) to fading out of existence in some manner worse than death. The windows didn't open wide enough for any three-dimensional person to squeeze through; the back door was right out, as it was locked and the various keys kept in our individual rooms. Besides, it was sticky and would jam every time we tried to open or close it. The various cupboards and cabinets were either too small or too full, and the small wooden table was both piled with empty takeaway containers and holding two wooden stools beneath it.

I heard myself trundling down the stairs inelegantly, and weakly staggered forward to find somewhere to hide. In doing so, I stepped back into the future, where the creature was once again before me, having finally accepted its situation somewhat and attempting to chew through the fabric covering itself. I half-turned, conflicted between addressing the unearthly creature before me or the impossible portal behind, but finally tiptoed sideways towards the creature's head, preferring the familiar risk of biting that came with any animal that I’d ever handled to another potential attack by its tail.

As a zoology student and animal rescue volunteer, I’d been trained on what to do in emergency rescue situations with feral or wild animals. The first and most ideal course of action was to focus on my own safety, and leave any particularly dangerous animals to the already trained professionals. But in my blind arrogance setting up this primitive trap, I’d already broken that first rule.

Lifting the gnawed remains of what had once been a fitted sheet for my mattress, I regretted my choice and drew back, eyes darting around in case I made any other strange travels. It wasn't just crocodilian: it had a long scaly body, and snapped powerful toothy jaws as I would expect from one, but its feet extended with long talons, and wide elephantine frills erupted from the back of its head as it attempted, not unsuccessfully, to intimidate me.

``It's okay,'' I told the thing. ``I'm not here to hurt you.''

Of course, the creature didn't understand me. It lunged forward, and I scrambled back until the wall of cupboards hit my back. The creature briefly paused, to the relief of my screaming nerves, before winding towards me in a serpentine manner despite its clawed feet. But before its pointed teeth could arrive close enough for me to flinch, cry, or defecate in petrification, it disappeared; I was left alone in the kitchen, staring blankly at thin air with only my aching body as proof that what had just transpired wasn't entirely a figment of my imagination. I reached forward to where the creature had vanished, cautiously prodding the air before me. Perhaps the creature was travelling to a time when nobody had been in the kitchen, for whatever strange reason. We rarely had meat in the house anyway, being a majority plant-based household. Was it here to forage?

When my hand didn't disappear and instead passed through the air as one would expect in a world where strange temporal passageways don't exist, I straightened and shakily strode to the sink to pour myself some water. It sloshed over the wavy brown rim as I held the mug, but I forcibly lifted it to my mouth, at which point an involuntary jerk of my wrist sent it sailing floorward to clatter into two separate pieces. I sighed, and picked up the remains before walking more slowly this time to the bin and pushing the foot-pedal.

The creature screeched and tore past me, bringing the contents of the trash along with it. In its mouth, it clutched a brown paper bag, and I inwardly cursed. Of course my one cheat meal a week from my mostly-vegan diet would be what attracted this strange creature. In the wild, you know not to bring any meat or fish that could attract predators. Now that this thing seemed to treat my kitchen as part of its territory, the same rule applied.

Watching its form slip past me from within the small bin, I had to marvel at its being able to fit inside; it was when I saw its head disappear into nothingness once more that I realised that it had never been entirely inside after all. And now it had littered my kitchen floor with day-old food remains! Against my better judgement and all of my training in animal management, I lunged forward and grabbed the reptile's tail, determined to follow it through its latest portal.

The creature immediately whipped its tail to flick me off, but I held fast. We passed through the gap into another kitchen, still my own but in the daytime and with a new mug upside down on the draining board. It had ``COFFEE'' written on it in a tall, skinny font a la Rae Dunn. Future me had some basic taste.

We continued to another version of my kitchen, and the thing snaked back towards me, snapping its monstrous teeth. I croaked rather than screamed, but before its jaws could clamp down on my skull, its tail whipped me away to its back, overpowering my grip and sending me flying into the cupboards once more.

``You really are stupid,'' I – another me – commented from where I sat at the small table, holding a large, misshapen hot pink mug that had to be some child's art project, perhaps inspired by a dragonfruit. ``Thanks to you, we're going to spend a fortune in mugs. Still are. Am.''

I didn't reply as I was still too winded to speak. In front of us, the reptile seemed not to realise that I was no longer attached to it and kept popping in and out of existence before us, making thundering rattles as it tried to shake off an imagined passenger for some time before it finally reappeared and lunged just in time to catch its own tail between its pointy teeth.

At the table, I took a sip and muttered, ``Unstoppable object meets immovable force.'' It sounded rehearsed.

Sitting in front of the creature, I finally got to my feet, though I couldn't bring myself to approach it again. ``What is it?'' I asked, intensely aware of the pain of being thrown across the kitchen twice now. ``And how can we get it out of the kitchen?''

``It's a wyrm. And why should I tell you? I didn't tell me,'' was my irritated reply.

I frowned, inspecting the reptile. Eyes, feet, claws... ``It's not a worm.''

``A wyrm,'' other me repeated. ``Like a dragon.''

``Dragons aren't real,'' I blurted automatically, already regretting it as the hideous monstrosity of a mug was slammed angrily onto the tabletop. ``Okay. It's a dragon.''

``A wyrm,'' I corrected myself, before turning away and pulling out my phone. ``Anyway, figure it out yourself.''

Well, that was unhelpful. Turning my attention away from myself, I watched the wyrm, which seemed to notice my gaze as it finally released its scaly tail to glare up at me.

``What do you want?'' I asked, but of course it didn't reply.

Instead, it charged toward me, and I didn't have the energy to dive away before it disappeared once again.

``Step forward,'' my bored voice instructed. I obeyed.

Nothing happened. Then I was hit with a force from behind which knocked the breath from me, and found myself travelling through variations of the kitchen at various times of day, all with myself watching from a distance with varying degrees of surprise.

``Wha–'' one asked, turning around from where I stood at the stove, dropping a floral pink mug before another me suddenly stood leaning against the counter, nodded at me and sipping from a speckled egg-blue teacup for a split-second until a third me waved from where I stood by the refrigerator with a pile of grey ceramic shards already at my feet. After several more visions of myself holding or dropping a variety of drinking vessels, I hit the floor, and the same self sat at the table once more, still scrolling and sipping from the same ugly mug. The wyrm now clambered over me, still skittering forward and leaving cuts and scratches on my skin and clothes as it passed over me.

``That hurt,'' I mumbled on the floor.

At this point, I was aching, miserable, and worst of all – tired. Even if the sun shone through the windows now, it had been the middle of the night when I had first confronted the wyrm. No intellectual curiosity could compel me at this point to understand the creature properly. All I wanted was to crawl up the stairs and into bed, and take a nice hot bath after approximately 300 years of snoozing.

``It hurt,'' other me agreed. “But did you work out why it’s like this?''

``Define like this,'' I muttered, closing my eyes slowly. Just for a little...

``It's anchored to this location, and it's extremely aggressive,'' my voice explained. I made a noise of acknowledgement. ``Wake up! Work out why it's being like this!''

With effort, I pried my eyelids apart and sat up.

``It's anchored to here? Is it hiding its offspring here?'' But it didn't seem like a very parental creature. Even if it had eggs, where would it find space to nest in somewhere as small as a houseshare kitchen?

``Or this place is just part of its territory?'' No, that didn't seem right. It had only nested here recently. Or maybe it had nested here a long time ago, just in the future. It would explain my general lack of surprise at the situation.

``Or is it unwell?'' That could be it. Why else would it remain somewhere as cloistered as my kitchen? It would prefer to remain in a small, cramped space to avoid any potential conspecifics or predators who would see it as an easy target. Predators... I shuddered at the thought of any beast large enough to see that monster as prey.

But if it were ill, how could I convince it to leave? Nurse it back to health? Forget spending money on mugs, I would have to spend all of my savings on meat if I were to. Given that the kitchen didn't seem full of carcasses and wyrm faeces, I doubted that the solution involved an extended period of nurture.

``Injured?'' I asked hesitantly, but I didn't reply, once again absorbed in my phone. That damn phone!

Watching myself slurp yet more coffee did make me wonder, though. How did I get through so many mugs? Surely I wouldn't have resorted to that ceramic atrocity unless forced to by circumstances. How many mugs was I doomed to break while plagued by the wyrm? Wouldn't the shards be a health risk? Shards!

I scrambled to my feet, and gasped as something sharp dug into my knee. Inspecting the fabric, I found my confirmation – a ceramic shard was buried in the fabric of my jeans, and it took some intense pinching to pull it out. Scanning the floor, I couldn't find any other scattered shards, but instead faint, gleaming gold prints. Crouching down to inspect one, a wild guess solidified to certainty.

``How can I trap it again?'' I queried, and the familiar, now very chewed mattress cover was tossed to me.

It would do. It wasn't optimal, but I didn't have any better options approaching. I stood back against the cabinets, wincing as the handles bumped against my sore ribs, and waited. When the wyrm reappeared this time, it finally stopped, and stood in place before me as it panted to repay its oxygen debt from the past few minutes of temporal cavorting. Resisting the desire to clamber onto the countertop and hide, I threw the fabric over the wyrm's face. The exhausted creature didn't respond, only lazily shaking its head before submitting after failing to dislodge the blind.

Slowly, cautiously, I crept around it and made what I hoped wouldn't be the most deadly mistake of my life. Given the creature's long bent legs and location on dry land in a kitchen less than a full two metres wide, I trusted that it wouldn't try – or succeed – in rolling me off. I perched on its back, preparing to be thrown into the next month, but fatigue and darkness seemed to make the wyrm docile. Grabbing one foot, I pulled it up off the floor to inspect the sole for any debris. Just like hoof-picking a donkey. Nothing was lodged in there, and the foot was free of any blood. The next two feet were the same, and the wyrm began to rumble its dissatisfaction as I finally lifted its fourth foot.

Crocodilians commonly known to exist on earth bleed, but very little. Vasoconstriction allows redirection of blood flow, and rapid blood clotting prevents any significant blood loss from the site of injury. The wyrm seemed the same, as the small ceramic splinter embedded in its food was surrounded by thin scabbing already, and even after I carefully removed it, no more than three thin, golden droplets of ichor-like serum dripped from within the wound before it began to form a translucent filmy scab where the intruding fragment had been.

The wyrm let out a loud rattle, and I shot to my feet, bolting to the door. To my relief, it didn't attempt any death rolls or anything of the sort, and only shook the mattress cover off its head before turning to me: the me at the door, staring down at the strange beast from where I had first observed it, wondering what strange creature was crashing about in my kitchen in the middle of the night. It turned and approached me, slowly, and my fingers gripped the door handle as its eyes seemed to lock onto mine. Gratitude? Anthropomorphism of animals was something that I usually avoided, but something in me hoped so. The wyrm stepped forward, made a small trilling noise, and then disappeared into a nothingness before me.

``Well done,'' I said from the table, as I sank to the floor in clammy relief.

My legs were jelly, my arms were jelly, my torso ached, and my head almost rang with exhaustion.

``What do I do now?'' I asked, or someone else seemed to with my mouth while my consciousness felt somewhere in the air above myself.\\
``You go home,'' I replied. ``Go to the countertop, and I'll help you through the wyrmhole it left for you.''\\
Muscles screaming, I obeyed, and stood uncertainly at the spot where I had arrived into this time.\\
``This will work?'' I asked, inexplicably distrusting of myself standing behind me with a malicious smile.\\
``Of course. I've been looking forward to this for months,'' I heard myself mutter, shoving myself forward onto the cold tiles of a blissfully empty kitchen.

It was over. My body ached. Rising to my feet, I finally made my way to the door to retreat to my bedroom and take a very long nap. As I took hold of the handle, there was a skittering noise, and I turned in time to see the wyrm vanish in its familiar panic. My future self's source of irritation was finally clear: the next few months were going to be very long.
